\documentclass[11pt]{article}
\usepackage{amsmath,amsthm,amssymb}
\usepackage[margin=1in]{geometry}

\newtheorem{theorem}{Theorem}
\newtheorem{observation}{Observation}
\newtheorem{conjecture}{Conjecture}
\theoremstyle{definition}
\newtheorem{example}{Example}
\theoremstyle{remark}
\newtheorem*{remark}{Remark}

\newcommand{\SYT}{\operatorname{SYT}}
\newcommand{\Des}{\operatorname{Des}}
\newcommand{\ord}{\operatorname{ord}}
\newcommand{\Z}{\mathbb{Z}}
\newcommand{\Q}{\mathbb{Q}}

\title{The graded vanishing-order law is a $d=2$ phenomenon:\\
collapse of $\ord_{\zeta_d} G_\lambda$ for $d \geq 3$}
\author{Clio}
\date{2026-06-06}

\begin{document}
\maketitle

\begin{abstract}
The graded refinement $G_\lambda(q) = \sum_{T} q^{s(T)}$ of the descent-counting
generating function carries, at $q=-1$, a rich vanishing-order law:
$\ord_{q=-1} G_\lambda = \lfloor |2\text{-core}(\lambda)|/2 \rfloor$, with order
growing without bound along the staircase of $2$-cores. It is natural to ask
whether this is the $d=2$ instance of a family of laws governing the vanishing of
$G_\lambda$ at primitive $d$-th roots of unity. The answer, by exhaustive
computation over all $271$ partitions of size $\le 12$ and all $d \in \{2,\dots,8\}$,
is no. For every $d \geq 3$ in this range the order $\ord_{\zeta_d} G_\lambda$ takes
only the values $0$ and $1$, and is positive exactly on the $d$-cores. The order law
does not lift; it stands alone at the involution $q=-1$. We explain why: the parity
weights $w_i$ are $0$ or odd, so $s(T)$ carries a clean mod-$2$ structure (the
domino / Murnaghan--Nakayama mechanism) and no analogous mod-$d$ structure for
$d \ge 3$.
\end{abstract}

\section{Setup}

Fix a partition $\lambda \vdash n$ and let $\SYT(\lambda)$ be the set of standard
Young tableaux of shape $\lambda$. For $T \in \SYT(\lambda)$ the \emph{descent set} is
\[
\Des(T) = \{\, i \in \{1,\dots,n-1\} : i+1 \text{ lies in a strictly lower row of } T
\text{ than } i \,\}.
\]
We weight the descents by the \emph{parity weight}
\[
w_i \;=\;
\begin{cases}
2i-1, & n-i \text{ odd},\\[2pt]
0, & n-i \text{ even},
\end{cases}
\qquad (1 \le i \le n-1),
\]
and form the \emph{descent statistic}
\[
s(T) \;=\; \sum_{i \in \Des(T)} w_i .
\]
The object of study is the \emph{graded polynomial}
\[
G_\lambda(q) \;=\; \sum_{T \in \SYT(\lambda)} q^{\,s(T)} \;\in\; \Z_{\ge 0}[q].
\]
This is the spectral grading of Theorem~B in the order-law program: the multiset
$\{s(T) : T \in \SYT(\lambda)\}$ is the multiset of branch exponents of the
Baxterised monodromy $\Omega$ on the irreducible $V^\lambda$, and $G_\lambda$ is its
generating function.

For a primitive $d$-th root of unity $\zeta_d$, write $\Phi_d(q)$ for the $d$-th
cyclotomic polynomial. Since $\Phi_d$ is the minimal polynomial of $\zeta_d$ over
$\Q$, the quantity
\[
\ord_{\zeta_d} G_\lambda \;=\; \max\{\, m \ge 0 : \Phi_d(q)^m \mid G_\lambda(q) \,\}
\]
is the order of vanishing of $G_\lambda$ at $\zeta_d$ --- the multiplicity with which
$G_\lambda$ vanishes as one approaches a primitive $d$-th root of unity. We study how
this order behaves as $d$ varies.

The single structural feature that drives everything below is this: \emph{every
nonzero weight $w_i = 2i-1$ is odd.} The weight set is two-valued in parity, taking
only the values $0$ and ``odd''. This is the crux.

\section{The $d=2$ law (recalled)}

At $d=2$ we have $\zeta_2 = -1$ and $\Phi_2(q) = q+1$. Here the program has a complete
theorem, proved earlier this cycle.

\begin{theorem}[Graded $2$-core vanishing law; Clio, 2026-06-01 / 06-02]
\label{thm:d2}
For every partition $\lambda$,
\[
\ord_{q=-1} G_\lambda \;=\; \left\lfloor \frac{|2\text{-core}(\lambda)|}{2} \right\rfloor,
\]
and the order depends only on the $2$-core of $\lambda$, not on its $2$-quotient.
Moreover the leading coefficient is explicit: writing $c$ for the order and $c_0$ for
the number of odd columns,
\[
G_\lambda(q) = a_c\,(q+1)^c + O\big((q+1)^{c+1}\big),
\qquad
a_c = \varepsilon(c_0)\,2^{-c}\, e_c(W_n)\,,
\]
with $W_n = \{\,2i-1 : n-i \text{ odd}\,\}$ and an explicit sign $\varepsilon$.
\end{theorem}

This is a \emph{rich} law. The $2$-cores are exactly the staircases
$\delta_k = (k,k-1,\dots,1)$, of sizes $0,1,3,6,10,15,\dots$, so the order is
unbounded over the family of all partitions. For instance, a partition whose $2$-core
is the staircase $(4,3,2,1)$ of size $10$ has $\ord_{q=-1} G_\lambda = 5$. The
vanishing at $q=-1$ can be as deep as one likes.

The mechanism is the domino / Murnaghan--Nakayama identity at $q=-1$: removing a
horizontal or vertical domino preserves the $2$-core, and the resulting
character-theoretic recursion forces $\chi^\lambda(2^a 1^b) = 0$ for $b$ below the
$2$-core threshold. The repeated factor $(q+1)$ is the polynomial shadow of this
$2$-adic rigidity.

\section{The main observation: collapse for $d \ge 3$}

One naturally hopes Theorem~\ref{thm:d2} is the $d=2$ member of a family --- a
``$\zeta_d$-lift'' in which $\lfloor |2\text{-core}|/2\rfloor$ is replaced by some
function of the $d$-core or $d$-quotient. We tested this directly.

\begin{remark}[Computational status]
The following is a \emph{computation}, not a theorem. For every $d \in \{2,\dots,8\}$
and every one of the $271$ partitions $\lambda$ with $|\lambda| \le 12$, the polynomial
$G_\lambda(q)$ was formed exactly and the multiplicity of $\Phi_d(q)$ dividing it was
computed over $\Q$. The results are tabulated below.
\end{remark}

\begin{center}
\renewcommand{\arraystretch}{1.25}
\begin{tabular}{c|c|c|l}
$d$ & behaviour & $\max_\lambda \ord_{\zeta_d}$ & support (where $\ord \ge 1$) \\
\hline
$2$ & \textbf{rich} & $5$ & $\ord = \lfloor |2\text{-core}|/2\rfloor$ exactly ($271/271$); nonzero on $59$ shapes\\
$3$ & degenerate & $1$ & only $\{(3,1),\,(2,1,1)\}$ --- the two $3$-cores of size $4$\\
$4$ & degenerate & $1$ & only $\{(2,2)\}$\\
$5$ & trivial & $0$ & $\ord \equiv 0$ \\
$6$ & degenerate & $1$ & $13$ shapes\\
$7$ & trivial & $0$ & $\ord \equiv 0$ \\
$8$ & trivial & $0$ & $\ord \equiv 0$ \\
\end{tabular}
\end{center}

The contrast with $d=2$ is sharp. Where the $q=-1$ order climbs the staircase without
bound, every $d \ge 3$ produces an order that never exceeds $1$, and that is positive
only on a thin set of exceptional shapes.

\begin{observation}[$d \ge 3$ collapse]
\label{obs:collapse}
For every $d \in \{3,\dots,8\}$ and every $\lambda$ with $|\lambda| \le 12$:
\[
\ord_{\zeta_d} G_\lambda \in \{0,1\},
\]
and $\ord_{\zeta_d} G_\lambda = 1$ only if $\lambda$ is itself a $d$-core (i.e.\ has
empty $d$-quotient). Equivalently, \emph{positive $d$-weight forces
$\ord_{\zeta_d} G_\lambda = 0$.} In particular there is no repeated cyclotomic factor:
$\Phi_d^2 \nmid G_\lambda$ throughout this range.
\end{observation}

Thus for $d \ge 3$ the only thing $\zeta_d$ can detect is a single, simple zero, and
it detects it only on $d$-cores. There is no growing order, hence nothing for a
``$\zeta_d$-order law'' to be \emph{about}.

\subsection*{Worked example: $\lambda = (3,1)$ at $d=3$}

Here $n = 4$, so the parity weights are
\[
(w_1,w_2,w_3) = \big(2\cdot 1 - 1,\; 0,\; 2\cdot 3 - 1\big) = (1,\,0,\,5),
\]
since $n - i = 4 - i$ is odd for $i \in \{1,3\}$ and even for $i = 2$. The shape
$(3,1)$ has three standard tableaux, distinguished by the position of the entry $4$
relative to its predecessor; their descent sets are $\{1\}$, $\{2\}$, $\{3\}$, giving
\[
s = w_1 = 1, \qquad s = w_2 = 0, \qquad s = w_3 = 5.
\]
Hence
\[
G_{(3,1)}(q) = q^{5} + q^{1} + q^{0} = q^5 + q + 1.
\]
This factors over $\Q$ as
\[
G_{(3,1)}(q) = (q^2 + q + 1)(q^3 - q^2 + 1) = \Phi_3(q)\cdot(q^3 - q^2 + 1),
\]
so $\ord_{\zeta_3} G_{(3,1)} = 1$: a single simple zero at $\zeta_3$, no more. The
cofactor $C(q) = q^3 - q^2 + 1$ evaluates at $\zeta_3 = e^{2\pi i/3}$ to
\[
C(\zeta_3) = \tfrac{5}{2} + \tfrac{\sqrt 3}{2}\, i,
\qquad |C(\zeta_3)|^2 = 7.
\]
This $7$ is the loud signal that the $d=2$ leading-coefficient structure does
\emph{not} lift: at $q=-1$ the leading coefficient is a manifestly factored
$2$-quotient hook product (Theorem~\ref{thm:d2}), but here the analogous cofactor has
squared modulus $7$ --- not a $3$-quotient hook product, not even a value with any
visible arithmetic meaning attached to the $3$-core. The fiber value is governed by
the classical Littlewood factorisation, but the \emph{order} is just $1$, and there is
no further structure underneath it.

\section{Why $d=2$ is special}

This section is explanation, not theorem. The phenomenon is forced by the parity
two-valuedness of the weights, and the explanation is correspondingly clean.

Every nonzero weight $w_i = 2i-1$ is odd. Therefore, at the involution $q = -1$,
\[
q^{s(T)} \big|_{q=-1} = (-1)^{s(T)},
\]
and $(-1)^{s(T)}$ depends only on the \emph{parity} of $s(T)$ --- equivalently, on the
parity of the number of \emph{odd-weight} descents of $T$. The statistic $s(T)$ thus
projects, at $q=-1$, onto a $\mathbb{Z}/2$-valued invariant. This mod-$2$ collapse is
precisely the domino structure: it is the polynomial avatar of the Murnaghan--Nakayama
recursion at $q=-1$, in which the relevant data is the $2$-core and the parity of
$2$-quotient statistics. The repeated factor $(q+1)$, and with it the entire
$\lfloor |2\text{-core}|/2\rfloor$ tower, is the visible trace of this $2$-adic
periodicity in the exponents.

For $d \ge 3$ there is simply no analogous mod-$d$ structure to exploit. The exponents
$s(T)$ are sums of values from $\{0\} \cup \{1,3,5,\dots\}$; reduced mod $d \ge 3$
these spread across the residues with no rigidity. A $d$-th root of unity $\zeta_d$
mixes the exponents according to their residue mod $d$, but no cancellation mechanism
forces \emph{repeated} vanishing, because there is no underlying mod-$d$ refinement of
the descent statistic. One zero --- the generic ``accidental'' one detected on
$d$-cores --- is all that can occur.

This is consistent with, and indeed predicted by, the earlier finding that
$G_\lambda(q)$ is \emph{generically irreducible over $\Q$} (2026-06-04). An
irreducible polynomial has no repeated factor at all, cyclotomic or otherwise; and
high-order vanishing at a root of unity is precisely the non-generic, structured
behaviour. The program established that the only place this non-generic behaviour
appears is the involution $q = -1$, driven by the $2$-core. The present computation is
the quantitative confirmation: at every other root of unity the order collapses to
$\{0,1\}$, exactly as a generically irreducible family should behave.

\section{Consequence for the program}

The upshot is a clean negative, and it sharpens the map of the territory.

\begin{itemize}
\item \textbf{No $\zeta_d$-order law exists for $d \ge 3$.} There is nothing to
characterise: with $\ord_{\zeta_d} G_\lambda \in \{0,1\}$ there is no
defect/obstruction term, no analogue of $\lfloor |2\text{-core}|/2\rfloor$, no
$d$-quotient correction to discover. The order is too shallow to carry information.

\item \textbf{Littlewood governs the value, not the order.} The classical
$d$-core / $d$-quotient factorisation does control the \emph{fiber} $G_\lambda(\zeta_d)$
--- the value of the polynomial at the root of unity --- and this is where the
factorisation lift genuinely lives (cf.\ the $2$-core entry of the program, where
$G_\lambda(\zeta_d)$ matches an Albion-type plethysm). But ``value at $\zeta_d$'' and
``order of vanishing at $\zeta_d$'' are different questions, and only the value lifts.

\item \textbf{The graded order law stands alone at $d=2$.} It is not the first member
of a family; it is a $2$-adic phenomenon with no higher-$d$ siblings.
\end{itemize}

Finally, the computation suggests a clean structural statement to aim for, which the
computation verifies but does not prove.

\begin{conjecture}
\label{conj:positive-weight}
For every $d \ge 3$ and every partition $\lambda$, if the $d$-weight of $\lambda$ is
positive (equivalently, the $d$-quotient is nonempty), then
\[
\ord_{\zeta_d} G_\lambda = 0.
\]
That is, $\Phi_d(q) \mid G_\lambda(q)$ forces $\lambda$ to be a $d$-core, and even
then the order is exactly $1$. (Verified for all $|\lambda| \le 12$ and
$d \in \{3,\dots,8\}$.)
\end{conjecture}

A proof would presumably proceed through the same fake-degree / cyclotomic-sieve
machinery that closed the $d=2$ law, but run at a $d$-th root of unity, where the
absence of mod-$d$ structure in $s(T)$ should manifest as the impossibility of a double
zero. That is the natural next target if one wishes to upgrade the present
observation to a theorem.

\end{document}
